\clearpage

\lehead[]{\normalfont\sffamily\hspace*{-2.00cm}\textcolor{white}{\colorbox{lightblue}{\makebox[1.60cm][r]{\thechapter}}}\hspace{0.17cm}\textcolor{lightblue}{\chaptertitle}}
\rohead[]{\textcolor{lightblue}{\chaptertitle}\normalfont\sffamily\hspace*{0.17cm}\textcolor{white}{\colorbox{lightblue}{\makebox[1.60cm][l]{\thechapter}}}\hspace{-2.00cm}}
%\chead[]{}
\rehead[]{\textcolor{lightblue}{AvHG, Inf, My}}
\lohead[]{\textcolor{lightblue}{AvHG, Inf, My}}

\section{Weiterführende \texttt{SELECT}-Techniken (Ausblick)}

SQL bietet noch eine Reihe weiterer nützlicher Möglichkeiten für Abfragen:

\begin{compactitem}
  \item \textbf{Ähnlichkeitssuche mit \lstinline|LIKE|:} Mit den Platzhaltern
  (\emph{Wildcards}) \lstinline|%| (beliebig viele, auch null Zeichen) \linebreak 
  und \lstinline|_| (genau ein beliebiges Zeichen) lassen sich Textmuster suchen,
  z.\,B.\ \linebreak
  \lstinline|SELECT * FROM schueler WHERE klasse LIKE '7%';|.
  \item \textbf{Berechnungen im \lstinline|SELECT|:} Man kann direkt in der
  Abfrage rechnen und dem Ergebnis einen Spaltennamen geben, z.\,B.\
  \lstinline|SELECT nummer, preis AS "Netto", (preis * 1.19) AS "Brutto" FROM artikel;|.
\end{compactitem}
